\section{Introduction}

\subsection{Problem statement and control objectives}

This report presents the modelling, simulation and experimental validation of control strategies for a magnetic levitation system. The objective of the experiment is to levitate a steel ball by controlling the magnetic force generated by an electromagnetic coil. The plant is challenging because the open-loop mechanical dynamics are unstable and the electromagnetic force is strongly nonlinear, since it depends on the square of the coil current and on the inverse square of the air gap.

For this reason, the control design is carried out around a selected equilibrium point and then validated both in simulation and on the real setup. The linearized model is used for the frequency-based and state-space controllers, while the final tracking and lift-up tasks require particular attention to saturation, operating range and nonlinear effects.

The work is organized around three main control objectives:

\begin{enumerate}[label=\textbf{O\arabic*.},leftmargin=1.4cm]
    \item \textbf{EM coil current control.} The electrical dynamics of the coil are regulated by an inner current loop. This loop is designed to be faster than the mechanical dynamics, so that the outer position controller can command a current reference instead of directly commanding the coil voltage.
    \item \textbf{Ball position stabilization.} The unstable mechanical dynamics are stabilized around a nominal levitation height. Frequency-based and state-space techniques are designed, simulated and compared.
    \item \textbf{Trajectory tracking and lift-up.} The system is tested outside the small-signal stabilization task, with reference trajectories and large-signal manoeuvres. In this phase, nonlinearities, actuator saturation and initialization become central issues.
\end{enumerate}
